%% 22_adctest.tex -- Chapter 22: Testing and Characterising ADCs
%% Status: OUTLINE (build stage 4). Replace the outline with the full chapter when it is written.
\chapter{Testing and Characterising ADCs}\label{ch:adctest}
\chapterlead{A converter is only as good as the measurement that proves it.}

\begin{specnote}[Chapter status: outline]
This chapter is an outline. It will be written in build stage 4 (the rest of Part III with Labs 7, 8 and 10). The topics, the design question and the labs below are fixed; the derivations, worked examples and figures are still to come.
\end{specnote}

\begin{objectives}
\begin{itemize}
  \item choose coherent sampling frequencies and windows for FFT tests;
  \item compute SNDR and SFDR from an FFT;
  \item measure DNL and INL with the sine-wave histogram and with static ramps;
  \item compare behavioural Verilog-A or Python models with transistor-level results;
  \item plan the measurements for a returned OpenMPW die.
\end{itemize}
\end{objectives}

\begin{insight}[Design question]
How do you prove, in simulation and on a returned die, that the \OB{} converters meet their specifications?
\end{insight}

\section{Coherent sampling and windowing}\label{sec:ch22-1}
Choosing $f_\mathrm{in} = (M/N)\fs$ with $M$ and $N$ co-prime; windows when coherence is impossible.

\section{FFT-based SNDR and SFDR}\label{sec:ch22-2}
Bins, leakage, harmonics and the noise floor; the processing gain of a long FFT.

\section{Histogram and ramp tests}\label{sec:ch22-3}
Code-density DNL and INL from a sine-wave histogram; static ramp tests for slow converters.

\section{Models versus transistors}\label{sec:ch22-4}
Behavioural models for fast system simulation and their validation against transistor-level results.

\section{Testing a returned die}\label{sec:ch22-5}
The bench set-up, the measurements that matter, and how to report them.

\begin{tryit}[Labs for this chapter]
Lab~6 (\cref{sec:lab6}) writes the \code{adctest.py} module used by Labs~8 and~10.
\end{tryit}

\paragraph{Read with this chapter} Pelgrom~\cite{pelgrom}.
